\documentclass[11pt,a4paper]{article}

\usepackage{kotex}
\usepackage{amsmath,amssymb,amsfonts,amsthm,bm}
\usepackage{graphicx}
\usepackage{pgfplots}
\pgfplotsset{compat=1.18}
\usepackage{booktabs}
\usepackage{algorithm}
\usepackage{algpseudocode}
\usepackage{float}
\usepackage{hyperref}
\usepackage{geometry}
\geometry{margin=25mm}

\title{Support Flow Tensor Field 기반 출력지향 메쉬 분할:\\
지지흐름 면특징을 이용한 클러스터링과 다방향 빌드 분할}
\author{설인환}
\date{}

\newtheorem{definition}{정의}
\newtheorem{theorem}{정리}
\newtheorem{proposition}{명제}

\begin{document}

\maketitle

\begin{abstract}
본 연구에서는 적층제조 출력을 고려한 메쉬 분할(mesh partitioning) 방법을 제안한다.
기존 분할 기법은 정점 좌표 $V\in\mathbb{R}^{N\times3}$ 또는 면 중심만을 입력으로 하여
DBSCAN, $k$-means, $k$-medoids 등으로 군집화하므로,
분할 경계에 ``출력 적합성''(오버행, 지지 구조, 빌드플레이트 접촉)이 전혀 반영되지 않는다.
본 연구는 Support Flow Tensor Field(SFTF)가 빌드 방향마다 면 단위로 계산하지만
3$\times$3 텐서와 스칼라 점수로 환원하면서 버리는 면별 지지흐름 물리량을 보존하여,
메쉬마다 자동 선택된 최적 빌드 방향에서 면별 특징행렬(adaptive mesh data)을 생성한다.
SFTF의 동반 연구에서 빌드플레이트 지지 비용마저 흐름 텐서에 가상 접지 노드를 더한
Rayleigh 항 $\tilde R=R+B$로 통합됨을 고려하면,
본 연구가 보존하는 면별 물리량은 바로 그 통합 흐름 텐서를 면별로 분해한 것에 해당한다.
이 특징행렬을 두 계열로 사용한다.
\textbf{(a) 특징 기반 군집화}는 SFTF 특징을 좌표와 융합하여
기존 $k$-medoids/DBSCAN/$k$-means에 입력한다.
\textbf{(b) 순수 지지흐름 분할}은 군집화 라이브러리 없이
SFTF의 지지 쌍과 역할만으로 메쉬를 나눈다
(지지 basin 분할 \texttt{support\_flow}, 영역 성장 \texttt{flow\_region}).
또한 상위 $K$개 빌드 방향에 대한 softmax--argmax 선호도로
각 부분이 서로 다른 좋은 빌드 방향을 갖도록 나누는 다방향 분할을 제시한다.
manikin($13{,}672$면)에서 기존 좌표 전용 $k$-means의 지지 클래스 순도는 $0.51$인 반면,
(a)와 (b)는 모두 $0.73$--$0.87$로 향상되어 분할 경계가 출력 지지 거동을 따름을 정량 확인하였다.
Stanford Bunny(69{,}662면)에서 면특징 추출은 약 $2.1$초,
각 분할은 $0.57$--$8.4$초에 수행되었고 모든 면이 빠짐없이 분할되었다.
\end{abstract}

\section{서론}

대형 또는 복잡한 형상을 적층제조로 출력할 때, 메쉬를 여러 부분으로 분할하여
각각을 개별 출력한 뒤 조립하는 전략이 널리 쓰인다.
이때 분할 경계는 단순히 기하학적으로 균등한 것이 아니라,
각 부분이 \emph{출력하기 좋은} 형태——즉 오버행과 지지구조가 적고
명확한 빌드 방향을 갖는 형태가 되도록 결정되어야 한다.

그러나 기존 메쉬 분할 파이프라인(예: \texttt{cut\_function})은
정점 좌표 $V$ 또는 면 중심 $F_c$만을 입력으로 하여
\begin{itemize}
\item DBSCAN: \texttt{DBSCAN(eps,min\_samples).fit(V)}
\item $k$-means: \texttt{k\_means(}$F_c$\texttt{,k)}
\item $k$-medoids: \texttt{KMedoids(k,metric="manhattan").fit(V)}
\item 응집형(agglomerative), 스펙트럴 군집화
\end{itemize}
와 같이 \emph{순수 좌표}에 대해 군집화한다.
이 경우 분할 기준은 공간적 근접성뿐이며,
``이 면이 특정 방향에서 오버행인가'', ``아래에 받쳐 줄 면이 있는가'',
``빌드플레이트까지 직접 지지되어야 하는가''와 같은 출력 정보는 반영되지 않는다.

한편 Support Flow Tensor Field(SFTF)~\cite{sftf}는 출력 방향 $n\in S^2$마다
면 단위로 오버행 계수, 면-면 지지 관계, 지지 높이, 빌드플레이트 접촉을 계산한다.
그러나 SFTF는 이 면별 물리량을 방향별 3$\times$3 텐서 $F(n)$과
스칼라 목적함수 $J(n)$으로 환원하여 \emph{최적 방향 랭킹}만을 출력하고,
면별 정보는 폐기한다.
최근 SFTF 개정에서는 이 환원을 더 밀고 나가, 빌드플레이트 penalty마저 흐름 텐서에
``법선이 빌드 방향인 가상 접지 노드''를 추가한 Rayleigh 항 $\tilde R=R+B$로 통합한다~\cite{sftf}.
즉 방향별 비용은 단일 통합 텐서의 스칼라 수축으로 수렴하는 반면,
그 텐서가 요약하는 면별 흐름장(flow field)——각 면이 어느 receiver로 흐르는지, 얼마나 높이 떨어져 있는지——은 폐기된다.
본 연구가 보존하려는 대상이 바로 이 면별 흐름장이다.
본 연구의 핵심 관찰은 다음과 같다.

\begin{quote}
군집화는 결국 $(N\times d)$ 특징행렬에 대해 수행된다.
기존 방법은 $d=3$(좌표)뿐이지만,
SFTF가 내부에서 이미 계산하는 면별 지지흐름 물리량을 행렬로 노출하면
그대로 군집화 입력으로 사용할 수 있다.
\end{quote}

따라서 본 연구는 SFTF를 \emph{메쉬 분할 구동기}로 재구성한다.
구체적 기여는 다음과 같다.
\begin{enumerate}
\item 메쉬별 최적 빌드 방향에서 면별 지지흐름 특징행렬을 생성하는
``adaptive mesh data'' 추출기(\S\ref{sec:features}).
\item 그 특징을 이용하는 세 가지 분할 경로:
지지영역 성장(\S\ref{sec:flowregion}),
좌표 융합 군집화(\S\ref{sec:cluster}),
다방향 softmax--argmax 빌드 분할(\S\ref{sec:multidir}).
\item 기존 점수 함수를 변경하지 않고 기존 SFTF 헬퍼를 재사용하는 구현과
예비 실험(\S\ref{sec:exp}).
\end{enumerate}

\section{면별 지지흐름 특징 추출}
\label{sec:features}

\subsection{빌드 방향의 적응적 선택}

분할의 기준이 되는 빌드 방향 $n^\ast$는 고정값이 아니라
SFTF 후보 평가에서 메쉬마다 자동으로 선택한다.
SFTF 후보 풀 $\{(J(n_k),\dots)\}$을 생성하고 각도 비최대억제(NMS)로
상위 방향을 고른 뒤, 최상위 방향을 $n^\ast$로 사용한다.
이 ``적응성''이 본 방법의 핵심으로,
면특징이 각 메쉬의 실제 최적 배향에 정렬된다.
다방향 확장(\S\ref{sec:multidir})에서는 상위 $K$개 방향을 모두 사용한다.

\subsection{면별 특징}

면 $i$의 중심을 $c_i$, 외향 단위법선을 $m_i$, 면적을 $A_i$,
빌드 방향을 $n=n^\ast$라 하자.
방향 $n$에서의 빌드플레이트 높이 기준은 $z_{\mathrm{plate}}(n)=\min_{v\in V}v\cdot n$이다.
각 면에 대해 다음 특징을 정의한다.

\begin{align}
O_i(n) &= \max(0,\,-m_i\cdot n) && \text{(오버행 계수)}\\
\tau_i(n) &= m_i\cdot n && \text{(부호 있는 기울기)}\\
\eta_i(n) &= c_i\cdot n - z_{\mathrm{plate}}(n) && \text{(빌드플레이트 위 높이)}
\end{align}

오버행 면($O_i>0$) 각각에 대해 SFTF와 동일하게 $c_i-\varepsilon n$에서 $-n$ 방향으로
ray를 발사하여 유효 receiver 면 $j$를 찾는다.
분할용 특징에서는 \emph{모든} 오버행 면에서 ray를 발사한다
(SFTF 채점에서 쓰는 $K_{\mathrm{ray}}$ stratified subsampling을 사용하지 않는다).
이로부터 다음을 면별로 기록한다.

\begin{itemize}
\item 지지 역할 $\rho_i\in\{\textsc{none},\textsc{bed},\textsc{face}\}$:
오버행이 아니면 \textsc{none},
유효 receiver를 찾으면 \textsc{face}(면-면 자기지지),
찾지 못하면 \textsc{bed}(빌드플레이트 직접 지지).
통합 텐서 관점에서 이 역할은 면 $i$의 지지 흐름이 향하는 receiver의 \emph{종류}와 같다.
\textsc{face}는 실제 면 receiver로의 흐름($m_im_j^T$ 항),
\textsc{bed}는 가상 접지 노드(법선 $n$)로의 흐름($m_in^T$ 항)이다~\cite{sftf}.
즉 $\rho_i$는 동반 연구의 증강 흐름 텐서를 면별로 분해한 것이며,
스칼라 환원에서 사라지는 ``어느 receiver로 흐르는가''라는 정보를 분할에 직접 사용한다.
\item 지지 높이 $h_i = h_{ij}(n)=(c_i-c_j)\cdot n$ (\textsc{face}일 때, 아니면 $0$).
\item 지지 대상 면 $t_i = j$ (없으면 $-1$).
\item 수신 횟수 $r_i$: 면 $i$를 향해 들어오는 지지 ray 수.
\end{itemize}

\subsection{특징행렬의 표준화}

면별 SFTF 특징을 열별로 $z$-점수 표준화한다.
값이 일정한 열은 $0$으로 둔다.
\begin{equation}
\Phi_{\mathrm{SFTF}}
= \operatorname{zscore}\!\Big(\big[\,O_i,\ \tau_i,\ \eta_i,\ h_i,\ r_i\,\big]\Big)
\in\mathbb{R}^{N\times 5}.
\end{equation}
표준화는 좌표(단위: 길이)와 지지량(단위: 무차원/길이/개수)의 척도 차이를 제거하여
거리 기반 군집화와 DBSCAN의 $\varepsilon$가 의미를 갖도록 하는 데 필수적이다.

\section{분할 방법}

세 가지 분할 경로는 모두 면별 라벨 $\ell\in\{0,\dots,K-1\}$(또는 잡음 $-1$)을 반환한다.

\subsection{순수 지지흐름 분할 (방법 b)}
\label{sec:flowregion}

군집화 라이브러리 없이 SFTF의 지지 물리량만으로 분할하는 두 변형을 제시한다.
이 계열은 좌표 군집화와 달리 sklearn 의존성이 없고,
분할 경계가 ``출력 지지 거동''을 직접 따른다.

\paragraph{지지흐름 basin 분할 (\texttt{support\_flow}).}
이 방법은 SFTF의 \emph{지지 쌍}(hit\_map)을 직접 사용한다.
먼저 각 면을 빌드 방향 $n^\ast$에서의 물리적 역할로 분류한다.
\begin{equation}
\kappa_i\in\{\textsc{bed},\ \textsc{needs},\ \textsc{supporter},\ \textsc{free}\},
\end{equation}
여기서 빌드플레이트 근처($\eta_i\le \delta\,\max_j\eta_j$)이면 \textsc{bed}(바닥접촉),
오버행이면 \textsc{needs}(지지 필요),
다른 면을 받치면($r_i>0$) \textsc{supporter}(지지함),
그 외는 \textsc{free}(자기지지)이다($\delta=0.02$ 기본값).
바닥접촉이 오버행보다 우선하는데, 이는 플레이트에 안착한 하향면은 지지가 불필요하기 때문이다.
이어 면-인접 그래프 위에서 다음 두 종류의 간선으로 union--find를 수행한다.
\begin{enumerate}
\item \emph{지지 쌍 간선}: 면지지 source $i$와 그 receiver $t_i$를 병합한다
(공간적으로 떨어진 패치라도 ``서로 받치는'' 면들을 하나로 묶는다).
\item \emph{동일 역할 인접 간선}: $\kappa_i=\kappa_j$인 인접면을 병합하여
각 역할 영역을 공간적으로 연결한다.
\end{enumerate}
결과적으로 각 부분은 ``오버행과 그것을 받치는 base''를 함께 포함하는
\emph{지지 basin}이 되어, 분할면을 따라 잘라도 지지 기둥이 한 부분 안에 유지된다.

\paragraph{기하적 지지영역 성장 (\texttt{flow\_region}).}
대안적으로, 지지 쌍 대신 법선 유사도로 영역을 성장시킨다.
인접한 두 면 $i,j$를
\begin{equation}
\rho_i=\rho_j \quad\text{이고}\quad m_i\cdot m_j \ge \cos\theta_{\mathrm{sim}}
\end{equation}
일 때에만 병합한다($\theta_{\mathrm{sim}}=35^\circ$ 기본값).
즉 같은 지지 역할이면서 법선이 충분히 유사한 면들을 한 영역으로 성장시킨다.
\texttt{support\_flow}가 지지 물리(쌍)를 강조한다면,
\texttt{flow\_region}은 표면 곡률 연속성을 강조하는 변형이다.

두 변형 모두 면 수가 $\tau_{\min}$ 미만인 작은 영역은
가장 많이 접한 이웃 영역으로 흡수하여 과분할을 억제한다.

\subsection{좌표 융합 군집화 (방법 a)}
\label{sec:cluster}

표준 군집화에 SFTF 특징을 추가 차원으로 주입한다.
표준화한 면 중심과 SFTF 특징을 가중 결합하여 군집화 행렬을 만든다.
\begin{equation}
M
= \big[\, w_s\,\operatorname{zscore}(C)\ \big|\ w_f\,\Phi_{\mathrm{SFTF}}\,\big]
\in\mathbb{R}^{N\times 8},
\label{eq:cluster-matrix}
\end{equation}
여기서 $C\in\mathbb{R}^{N\times3}$는 면 중심, $w_s,w_f$는 좌표/특징 가중치이다.
이후 $M$에 대해 다음 중 하나를 적용한다.
\begin{itemize}
\item \textbf{DBSCAN/$k$-means/agglomerative}: scikit-learn 구현을 $M$에 적합.
\item \textbf{$k$-medoids}: 대형 메쉬에서 $O(N^2)$ 거리행렬을 피하기 위해
교대 갱신형 확장 $k$-medoids를 자체 구현한다.
각 반복은 (i) 현 medoid까지의 거리로 $O(NK)$ 군집 배정,
(ii) 군집 평균에 가장 가까운 실제 구성원으로 medoid를 갱신한다.
거리는 manhattan을 사용하며 난수 시드를 고정하여 결정론적이다.
\end{itemize}
DBSCAN의 잡음 면(라벨 $-1$)은 분할 부분에서 제외할 수 있다.

\subsection{다방향 softmax--argmax 빌드 분할 (방법 a+b)}
\label{sec:multidir}

단일 방향만 사용하면 모든 부분이 같은 빌드 방향을 전제한다.
각 부분이 \emph{서로 다른} 좋은 빌드 방향을 갖게 하려면,
상위 $K$개 SFTF 방향 $\{n_1,\dots,n_K\}$ 각각에 대한 면별 비용을 쌓는다.
면 $i$의 방향 $n_k$에 대한 지지 비용을
\begin{equation}
c_{ik} = O_i(n_k) = \max(0,\,-m_i\cdot n_k)
\end{equation}
로 두면(저비용, ray 불필요), 비용행렬 $C\in\mathbb{R}^{N\times K}$를 얻는다.
선택적으로 ``support'' 모드에서는 방향마다 전체 면특징을 추출하여
지지기둥 높이와 빌드플레이트 접촉까지 비용에 반영한다($K$회 ray casting).

면별 선호 방향과 soft 가중치를
\begin{align}
b_i &= \arg\min_k c_{ik}
   && \text{(그 면을 가장 덜 오버행으로 만드는 방향)}\\
w_{ik} &= \operatorname{softmax}_k(-\beta\, c_{ik})
   = \frac{e^{-\beta c_{ik}}}{\sum_{k'} e^{-\beta c_{ik'}}}
   && (\textstyle\sum_k w_{ik}=1)
\end{align}
로 정의한다($\beta=8$ 기본값).
$b_i$를 면별 ``regime''으로 사용하여 \S\ref{sec:flowregion}의 영역 성장을 수행하면,
각 부분은 \emph{단일하고 일관된} 빌드 방향을 갖는 연결 영역이 된다.
한편 soft 가중치 $w_{ik}$는 식~\eqref{eq:cluster-matrix}에
$w_d\,[\,w_{i1},\dots,w_{iK}\,]$ 블록으로 덧붙여
$k$-medoids/DBSCAN에도 방향 선호를 융합할 수 있다(가중치 $w_d$).

\section{알고리즘과 구현}

\begin{algorithm}[H]
\caption{출력지향 메쉬 분할}
\begin{algorithmic}[1]
\State 후보 풀에서 상위 빌드 방향 $n^\ast$(또는 상위 $K$개 $\{n_k\}$)를 선택
\For{each face $i$}
  \State $O_i\gets\max(0,-m_i\cdot n^\ast)$, $\tau_i\gets m_i\cdot n^\ast$,
         $\eta_i\gets c_i\cdot n^\ast-z_{\mathrm{plate}}$
\EndFor
\For{each overhang face $i$ ($O_i>0$)}
  \State $c_i-\varepsilon n^\ast$에서 $-n^\ast$로 ray 발사, 유효 receiver $j$ 탐색
  \If{$j$ 존재} $\rho_i\gets\textsc{face}$, $h_i\gets(c_i-c_j)\cdot n^\ast$, $r_j\gets r_j+1$
  \Else\ $\rho_i\gets\textsc{bed}$ \EndIf
\EndFor
\State $\Phi_{\mathrm{SFTF}}\gets\operatorname{zscore}([O,\tau,\eta,h,r])$
\If{method = support\_flow}\Comment{방법 b: 지지 물리}
  \State 각 면을 $\kappa_i\in\{$bed, needs, supporter, free$\}$로 분류
  \State union--find: 지지 쌍 $(i,t_i)$ 병합 $\wedge$ $\kappa_i{=}\kappa_j$ 인접면 병합
  \State 작은 영역 흡수(선택)
\ElsIf{method = flow\_region 또는 build\_direction}\Comment{방법 b: 기하/방향}
  \State regime $\gets \rho$ (또는 다방향 $b$)
  \State $\ell\gets$ 면-인접 union--find: regime 동일 $\wedge$ $m_i\cdot m_j\ge\cos\theta_{\mathrm{sim}}$일 때 병합
  \State 작은 영역 흡수(선택)
\Else\Comment{방법 a: 특징 군집화}
  \State $M\gets[w_s\,\operatorname{zscore}(C)\mid w_f\,\Phi_{\mathrm{SFTF}}\mid w_d\,W]$
  \State $\ell\gets$ \{DBSCAN, $k$-medoids, $k$-means, agglomerative\} 중 적용
\EndIf
\State \Return 면별 라벨 $\ell$; 라벨별 submesh로 분리
\end{algorithmic}
\end{algorithm}

구현은 기존 SFTF 모듈의 헬퍼(표면 데이터 캐시, ray sampling, 유효 hit 판정)를
\emph{변경 없이} 재사용하며, 테스트로 보호된 방향별 점수 함수
$F(n),J(n)$은 수정하지 않는다.
면특징 추출은 별도 경로로 동일한 ray casting 원천 함수를 호출한다.
복잡도는 면특징 추출이 ray 수에 선형,
영역 성장이 인접 모서리 수에 거의 선형($O(E\,\alpha(N))$),
확장 $k$-medoids가 반복당 $O(NK)$이다.

\section{예비 실험}
\label{sec:exp}

Stanford Bunny($N=69{,}662$면)에 대해 단일 추출을 공유하여 세 방법을 적용하였다.
면특징 추출은 약 $2.1$초였고,
자동 선택된 빌드 방향은 $n^\ast\approx(0.78,0.58,0.24)$,
오버행 면은 $35{,}933/69{,}662$였다.

\begin{table}[H]
\centering
\caption{Bunny(69{,}662면) 분할 결과.
모든 비잡음 방법에서 분할 부분의 면 수 합은 원본과 동일하다(면 보존).}
\label{tab:partition}
\begin{tabular}{lrrr}
\toprule
방법 & 시간(s) & 부분 수 & 잡음 면\\
\midrule
support\_flow (방법 b, 지지 basin)   & 0.69 & 4  & 0\\
flow\_region (방법 b, 영역 성장)     & 0.57 & 9  & 0\\
$k$-medoids (방법 a, 좌표+특징 융합)  & 0.68 & 6  & 0\\
DBSCAN (방법 a, 좌표+특징 융합)      & 8.36 & 16 & 347\\
build\_direction (다방향 argmax)    & 0.59 & 5  & 0\\
\bottomrule
\end{tabular}
\end{table}

다방향 분할에서는 상위 $K=3$ 방향에 대해 면이
방향 0에 $51.8\%$, 방향 1에 $48.2\%$로 갈렸으며
(세 번째 방향은 첫 방향과 거의 같은 축이어서 흡수됨),
영역 성장 결과 5개 연결 부분이 두 빌드 방향(3개$\to$방향1, 2개$\to$방향0)에 사상되었다.
모든 방법에서 전 면이 분할되었다.

\subsection{기존 좌표 기반 군집화와의 정량 비교}
\label{sec:purity}

Figure~\ref{fig:partition-capture}는 manikin에 다섯 분할 방법을 적용한 실제 결과로,
\texttt{run\_mesh\_partition.py}의 직교(평행) 투영 렌더링을 캡처한 것이다.
면색은 부분 라벨을 나타낸다.
순수 지지흐름 방법(\texttt{support\_flow}, \texttt{flow\_region})과
특징 융합 방법($k$-medoids, DBSCAN)은 분할 경계가 오버행/지지 영역 구조를 따르는 반면,
다방향 \texttt{build\_direction}은 빌드 방향 선호가 균질한(73\%가 동일 방향) manikin에서
대부분 단일 영역으로 묶인다.

\begin{figure}[H]
\centering
\begin{tikzpicture}
\begin{axis}[
  width=\linewidth,
  axis equal image,
  enlargelimits=false,
  axis on top,
  xmin=0, xmax=1, ymin=0, ymax=0.3158,
  axis y line=none,
  axis x line=bottom,
  x axis line style={draw=none},
  major tick length=0pt,
  xtick={0.1183,0.3091,0.5000,0.6909,0.8817},
  xticklabels={\texttt{support\_flow},\texttt{flow\_region},$k$-medoids,\textsc{dbscan},\texttt{build\_direction}},
  xticklabel style={font=\footnotesize},
]
\addplot graphics [xmin=0,xmax=1,ymin=0,ymax=0.3158] {sftf_partition_manikin.png};
\end{axis}
\end{tikzpicture}
\caption{manikin($13{,}672$면)에 대한 다섯 분할 방법의 적용 결과(직교 투영 렌더링 캡처).
면색은 부분 라벨이며, 좌측부터 순수 지지흐름(방법 b) 두 가지, 특징 융합(방법 a) 두 가지,
다방향 빌드 분할이다. 이미지는 pgfplots의 그래픽 삽입 기능(\texttt{addplot graphics})으로 넣고
열별 방법명을 축 눈금으로 주석하였다.}
\label{fig:partition-capture}
\end{figure}

본 방법의 핵심 주장——``분할 경계가 출력 지지 거동을 따른다''——을 정량화하기 위해
\emph{지지 클래스 순도}(support-class purity)를 정의한다.
면별 지지 클래스 $\kappa_i\in\{$bed, needs, supporter, free$\}$에 대해,
분할 $\{\mathcal{P}_\ell\}$의 순도는 각 부분에서 지배 클래스가 차지하는 비율의
면 수 가중 평균이다.
\begin{equation}
\mathrm{Purity}
= \frac{1}{N}\sum_\ell \max_{\kappa}
\big|\{i\in\mathcal{P}_\ell : \kappa_i=\kappa\}\big|.
\end{equation}
순도가 $1$에 가까울수록 각 부분이 단일 지지 거동(예: 오버행만, 또는 자기지지만)으로
구성되어 출력 적합성이 높다.

Table~\ref{tab:purity}는 manikin($13{,}672$면)에 대해
기존 좌표 전용 $k$-means와 본 연구의 (a)/(b) 방법을 비교한다.
좌표만 사용하는 기존 $k$-means는 순도 $0.51$로,
한 부분 안에 오버행과 자기지지 면이 섞인다.
SFTF 특징을 융합한 방법 (a)와 순수 지지흐름 방법 (b)는 모두 순도 $0.73$--$0.87$로
크게 향상되어, 분할 경계가 실제 지지 영역 구조를 따름을 정량적으로 확인한다.

\begin{table}[H]
\centering
\caption{manikin($13{,}672$면)에 대한 분할 방법별 지지 클래스 순도 비교(높을수록 좋음).
기준선은 좌표만 사용하는 기존 $k$-means이다.}
\label{tab:purity}
\begin{tabular}{llrr}
\toprule
계열 & 방법 & 부분 수 & 순도\\
\midrule
기준선 & $k$-means (좌표 전용)          & 6  & 0.508\\
\midrule
(a)   & $k$-means (좌표+SFTF 특징)     & 6  & \textbf{0.873}\\
(a)   & $k$-medoids (좌표+SFTF 특징)   & 6  & 0.745\\
(a)   & DBSCAN (좌표+SFTF 특징)        & 13 & 0.566\\
\midrule
(b)   & flow\_region (영역 성장)       & 29 & 0.725\\
(b)   & support\_flow (지지 basin)     & 16 & 0.749\\
\bottomrule
\end{tabular}
\end{table}

기존 좌표 기반 군집화는 공간적 근접성만 보므로 출력 적합성과 무관하게 분할되는 반면,
(a)는 동일한 군집화 알고리즘에 SFTF 특징을 추가하는 것만으로 순도를 $0.51\to0.87$로 끌어올리고,
(b)는 군집화 라이브러리 없이도 지지 물리만으로 $0.73$ 이상의 순도를 달성한다.
모든 방법은 manikin에서 $1$초 이내(특징 추출 별도)에 수행되었다.

\section{결론 및 향후 연구}

본 연구는 SFTF의 면별 지지흐름 물리량을 보존하여
출력 적합성을 직접 반영하는 메쉬 분할 방법을 제안하였다.
두 계열——(a) SFTF 특징을 기존 군집화에 융합하는 방법과
(b) 지지 쌍/역할만으로 분할하는 순수 지지흐름 방법——은
모두 좌표 전용 군집화 대비 지지 클래스 순도를 크게 높였으며,
다방향 softmax--argmax 분할은 각 부분이 서로 다른 좋은 빌드 방향을 갖도록 한다.
방법론적으로 본 연구는 동반 연구가 면별 흐름장을 단일 통합 텐서의 Rayleigh 수축
$\tilde R=R+B$로 환원하는 것과 정확히 상보적이다~\cite{sftf}.
동반 연구는 그 통합 텐서로 \emph{방향}을 고르고, 본 연구는 같은 흐름장의 면별 분해(특히 receiver 종류 $\rho_i$)로 \emph{분할}을 수행한다.

향후 연구로는 (i) 각 부분의 빌드 방향을 TOMO\_INT3 또는 슬라이서로 재검증하여
분할 품질을 정량화하고,
(ii) ``support'' 비용 모드와 후보 방향의 직교성 제어로 다방향 분할의 방향 다양성을 높이며,
(iii) 절단면 면적/접합 강도/조립성 항을 더한 다목적 분할로 확장할 계획이다.

\nocite{dbscan,kmedoids}
\bibliographystyle{plain}
\bibliography{references}

\end{document}
